\documentclass[10pt,a4paper]{amsart}
\usepackage{amsmath,amssymb,amsthm,graphicx}
\usepackage{verbatim}
\newenvironment{reference}{\comment}{\endcomment}
%\newenvironment{reference}{}{}
\newenvironment{slogan}{\comment}{\endcomment}
\newenvironment{history}{\comment}{\endcomment}
% For commutative diagrams we use Xy-pic
\usepackage[all]{xy}

% We use 2cell for 2-commutative diagrams.
\xyoption{2cell}
% Theorem environments.
%
\theoremstyle{plain}
\newtheorem{theorem}[subsection]{Theorem}
\newtheorem{proposition}[subsection]{Proposition}
\newtheorem{lemma}[subsection]{Lemma}

\theoremstyle{definition}
\newtheorem{definition}[subsection]{Definition}
\newtheorem{example}[subsection]{Example}
\newtheorem{exercise}[subsection]{Exercise}
\newtheorem{situation}[subsection]{Situation}

\theoremstyle{remark}
\newtheorem{remark}[subsection]{Remark}
\newtheorem{remarks}[subsection]{Remarks}

\numberwithin{equation}{subsection}

% Macros
%
\def\lim{\mathop{\mathrm{lim}}\nolimits}
\def\colim{\mathop{\mathrm{colim}}\nolimits}
\def\Spec{\mathop{\mathrm{Spec}}}
\def\Hom{\mathop{\mathrm{Hom}}\nolimits}
\def\Ext{\mathop{\mathrm{Ext}}\nolimits}
\def\SheafHom{\mathop{\mathcal{H}\!\mathit{om}}\nolimits}
\def\SheafExt{\mathop{\mathcal{E}\!\mathit{xt}}\nolimits}
\def\Sch{\mathit{Sch}}
\def\Mor{\mathop{\mathrm{Mor}}\nolimits}
\def\Ob{\mathop{\mathrm{Ob}}\nolimits}
\def\Sh{\mathop{\mathit{Sh}}\nolimits}
\def\NL{\mathop{N\!L}\nolimits}
\def\CH{\mathop{\mathrm{CH}}\nolimits}
\def\proetale{{pro\text{-}\acute{e}tale}}
\def\etale{{\acute{e}tale}}
\def\QCoh{\mathit{QCoh}}
\def\Ker{\mathop{\mathrm{Ker}}}
\def\Im{\mathop{\mathrm{Im}}}
\def\Coker{\mathop{\mathrm{Coker}}}
\def\Coim{\mathop{\mathrm{Coim}}}

% Boxtimes
%
\DeclareMathSymbol{\boxtimes}{\mathbin}{AMSa}{"02}

%
% Macros for moduli stacks/spaces
%
\def\QCohstack{\mathcal{QC}\!\mathit{oh}}
\def\Cohstack{\mathcal{C}\!\mathit{oh}}
\def\Spacesstack{\mathcal{S}\!\mathit{paces}}
\def\Quotfunctor{\mathrm{Quot}}
\def\Hilbfunctor{\mathrm{Hilb}}
\def\Curvesstack{\mathcal{C}\!\mathit{urves}}
\def\Polarizedstack{\mathcal{P}\!\mathit{olarized}}
\def\Complexesstack{\mathcal{C}\!\mathit{omplexes}}
% \Pic is the operator that assigns to X its picard group, usage \Pic(X)
% \Picardstack_{X/B} denotes the Picard stack of X over B
% \Picardfunctor_{X/B} denotes the Picard functor of X over B
\def\Pic{\mathop{\mathrm{Pic}}\nolimits}
\def\Picardstack{\mathcal{P}\!\mathit{ic}}
\def\Picardfunctor{\mathrm{Pic}}
\def\Deformationcategory{\mathcal{D}\!\mathit{ef}}
\usepackage{cleveref}
\crefname{originalreference}{original reference}{original references}
\Crefname{originalreference}{Original reference}{Original references}
\makeatletter
\newlabel{algebra-lemma-geometrically-integral-any-integral-base-change}{{algebra-lemma-geometrically-integral-any-integral-base-change}{source}{}{sourcepointer-1}{}}
\newlabel{algebra-lemma-geometrically-integral-any-integral-base-change@cref}{{[originalreference][0][]algebra-lemma-geometrically-integral-any-integral-base-change}{[1][0][]source}{}{}{}}
\makeatother
\crefname{originallemma}{original lemma}{original lemmas}
\Crefname{originallemma}{Original lemma}{Original lemmas}
\makeatletter
\newlabel{lemma-universally-injective}{{lemma-universally-injective}{source}{}{sourcepointer-2}{}}
\newlabel{lemma-universally-injective@cref}{{[originallemma][0][]lemma-universally-injective}{[1][0][]source}{}{}{}}
\makeatother
\crefname{originalreference}{original reference}{original references}
\Crefname{originalreference}{Original reference}{Original references}
\makeatletter
\newlabel{stacks-properties-lemma-unique-point}{{stacks-properties-lemma-unique-point}{source}{}{sourcepointer-3}{}}
\newlabel{stacks-properties-lemma-unique-point@cref}{{[originalreference][0][]stacks-properties-lemma-unique-point}{[1][0][]source}{}{}{}}
\makeatother
\crefname{originalreference}{original reference}{original references}
\Crefname{originalreference}{Original reference}{Original references}
\makeatletter
\newlabel{varieties-proposition-unique-base-field}{{varieties-proposition-unique-base-field}{source}{}{sourcepointer-4}{}}
\newlabel{varieties-proposition-unique-base-field@cref}{{[originalreference][0][]varieties-proposition-unique-base-field}{[1][0][]source}{}{}{}}
\makeatother
\title{A point of the target of a universally injective morphism over an algebraically closed field lifts to a point of the source}
\author{The Stacks Project Authors}
\date{stacks-project:stacks-morphisms.tex:local-snapshot-20261012; GNU FDL 1.2 or later}
\begin{document}
\maketitle
\noindent\textit{Editorial note.} All delivered mathematical text is the author's own, copied from the pinned Stacks Project chapter source sources/stacks-algebraic-stacks-morphisms/pinned-original-source.tex; no statement or proof step was rewritten, reordered or removed.\par\medskip
\noindent\textit{Editorial note.} Preamble provenance: the Stacks Project keeps one single shared preamble for every chapter; this card takes it directly from the staged sources/stacks-algebraic-stacks-morphisms/preamble.tex as three verbatim spans (lines 8-12, 14-18 and 154-217), so no derived file was written and nothing was added or redefined. The wrapper class and the base packages are supplied by the assembler.\par\medskip
\noindent\textit{Editorial note.} Bibliography provenance: neither the counted statement nor its proof contains a citation (this lane rescans the whole span before writing the card), so the card carries no bibliography and no entry was reconstructed.\par\medskip
\noindent\textit{Editorial note.} Selection provenance: this card was opened under the file-based control contract (CONTROL.md), with the mechanical pre-card gate run on its statement and proof interval before the card was written.\par\medskip
\noindent\textit{Allowed standard prerequisites (editorial).}\begin{itemize}
\item The existence of a surjective flat locally of finite type cover of a nonempty algebraic stack by the spectrum of a field. Invoked inside the counted proof (source lines 2997-2998) as a cross-chapter pointer. Produces the field extension through which the reduction is carried out.
\item The characterization of universal injectivity through surjectivity of the diagonal. Invoked inside the counted proof (source lines 2999-3002). Makes the relevant morphism of fibre products surjective.
\item The statement that a tensor product of a field with itself over an algebraically closed base is a domain. Invoked inside the counted proof (source lines 3003-3006) as a cross-chapter pointer. Provides the generic point used to compare the two projections.
\item The statement on equal integral closures of the images of a field under two morphisms. Invoked near the end of the counted proof (source lines 3013-3016) as a cross-chapter pointer. Forces the comparison of the two images to give the desired equality of fields.
\end{itemize}\medskip
\ifdefined\hypertarget\hypertarget{sourcepointer-1}{}\fi
\noindent\textit{Source reference.} Original reference algebra-lemma-geometrically-integral-any-integral-base-change: sources/stacks-algebraic-stacks-morphisms/pinned-original-source.tex, source lines 3007--3007. Navigation pointer only: the counted text refers to the Stacks item algebra-lemma-geometrically-integral-any-integral-base-change, which belongs to another chapter of the project and is not part of this card. The pointer carries the target's own label as its number because this chapter's staged snapshot does not contain the other chapter's tag table; its content is not reproduced here.\par\medskip
\ifdefined\hypertarget\hypertarget{sourcepointer-2}{}\fi
\noindent\textit{Source reference.} Original lemma lemma-universally-injective: sources/stacks-algebraic-stacks-morphisms/pinned-original-source.tex, source lines 2907--2907. Navigation pointer only: the counted text refers to the Stacks item lemma-universally-injective, which is an existing labeled result of this same chapter and is not part of this card. The pointer carries the target's own label as its number and its definition line as its span; its content is not reproduced here.\par\medskip
\ifdefined\hypertarget\hypertarget{sourcepointer-3}{}\fi
\noindent\textit{Source reference.} Original reference stacks-properties-lemma-unique-point: sources/stacks-algebraic-stacks-morphisms/pinned-original-source.tex, source lines 2997--2997. Navigation pointer only: the counted text refers to the Stacks item stacks-properties-lemma-unique-point, which belongs to another chapter of the project and is not part of this card. The pointer carries the target's own label as its number because this chapter's staged snapshot does not contain the other chapter's tag table; its content is not reproduced here.\par\medskip
\ifdefined\hypertarget\hypertarget{sourcepointer-4}{}\fi
\noindent\textit{Source reference.} Original reference varieties-proposition-unique-base-field: sources/stacks-algebraic-stacks-morphisms/pinned-original-source.tex, source lines 3017--3017. Navigation pointer only: the counted text refers to the Stacks item varieties-proposition-unique-base-field, which belongs to another chapter of the project and is not part of this card. The pointer carries the target's own label as its number because this chapter's staged snapshot does not contain the other chapter's tag table; its content is not reproduced here.\par\medskip
\begin{lemma}
\label{lemma-universally-injective-point}
Let $f : \mathcal{X} \to \mathcal{Y}$ be a universally injective
morphism of algebraic stacks. Let $y : \Spec(k) \to \mathcal{Y}$
be a morphism where $k$ is an algebraically closed field.
If $y$ is in the image of $|\mathcal{X}| \to |\mathcal{Y}|$,
then there is a morphism $x : \Spec(k) \to \mathcal{X}$
with $y = f \circ x$.
\end{lemma}

\begin{proof}
We first remark this lemma is not a triviality, because the assumption that
$y$ is in the image of $|f|$ means only that we can lift
$y$ to a morphism into $\mathcal{X}$ after possibly replacing
$k$ by an extension field. To prove the lemma we may base change
$f$ by $y$, hence we may assume we have a nonempty algebraic stack
$\mathcal{X}$ and a universally injective morphism
$\mathcal{X} \to \Spec(k)$ and we want to find a $k$-valued point
of $\mathcal{X}$. We may replace $\mathcal{X}$ by its reduction.
We may choose a field $k'$ and a surjective, flat, locally finite type morphism
$\Spec(k') \to \mathcal{X}$, see
Properties of Stacks, Lemma \ref{stacks-properties-lemma-unique-point}.
Since $\mathcal{X} \to \Spec(k)$ is universally injective, we find that
$$
\Spec(k') \times_\mathcal{X} \Spec(k') \to \Spec(k' \otimes_k k')
$$
is surjective as the base change of the surjective morphism
$\Delta : \mathcal{X} \to \mathcal{X} \times_{\Spec(k)} \mathcal{X}$
(Lemma \ref{lemma-universally-injective}).
Since $k$ is algebraically closed $k' \otimes_k k'$ is a domain
(Algebra, Lemma
\ref{algebra-lemma-geometrically-integral-any-integral-base-change}).
Let $\xi \in \Spec(k') \times_\mathcal{X} \Spec(k')$
be a point mapping to the generic point of $\Spec(k' \otimes_k k')$.
Let $U$ be the reduced induced closed subscheme structure on
the connected component of $\Spec(k') \times_\mathcal{X} \Spec(k')$
containing $\xi$. Then the two projections $U \to \Spec(k')$
are locally of finite type, as this was true for the projections
$\Spec(k') \times_\mathcal{X} \Spec(k') \to \Spec(k')$
as base changes of the morphism $\Spec(k') \to \mathcal{X}$.
Applying
Varieties, Proposition \ref{varieties-proposition-unique-base-field}
we find that the integral closures of the two images
of $k'$ in $\Gamma(U, \mathcal{O}_U)$ are equal.
Looking in $\kappa(\xi)$ means that any element of the form
$\lambda \otimes 1$ is algebraically dependent on
the subfield
$$
1 \otimes k' \subset 
(\text{fraction field of }k' \otimes_k k') \subset
\kappa(\xi).
$$
Since $k$ is algebraically closed, this is only possible
if $k' = k$ and the proof is complete.
\end{proof}

\end{document}
